\documentclass[11pt,a4paper]{article}

\usepackage[margin=2.5cm]{geometry}
\usepackage{booktabs}
\usepackage{graphicx}
\usepackage{amsmath}
\usepackage{hyperref}
\usepackage{float}
\usepackage{multirow}
\usepackage{array}
\usepackage{parskip}
\usepackage{caption}
\usepackage{subcaption}
\usepackage{cite}
\usepackage{xcolor}
\usepackage{siunitx}
\sisetup{detect-weight=true, detect-inline-weight=math}

\hypersetup{colorlinks=true, linkcolor=blue, citecolor=blue, urlcolor=blue}

\title{\textbf{HC701 -- Assignment 3}\\
Ultrasound Nerve Segmentation}
\author{Mashrafi Monon}
\date{April 2026}

\begin{document}

\maketitle
\tableofcontents
\newpage

%% ─────────────────────────────────────────────────────────────────────────────
\section{Task 1.1 -- Experimental Design and Justification}
%% ─────────────────────────────────────────────────────────────────────────────

\subsection{Dataset and Preprocessing}

The dataset consists of 2,757 grayscale ultrasound images of the brachial plexus nerve region
(580$\times$420 pixels) with corresponding binary segmentation masks, and a separate test set of
2,878 images.
All images were resized to 256$\times$256 pixels and normalised using ImageNet mean
and standard deviation (mean = [0.485, 0.456, 0.406], std = [0.229, 0.224, 0.225]).
Grayscale images were replicated to three channels to be compatible with ImageNet-pretrained
encoders, a common and well-validated practice in medical image segmentation
\cite{ronneberger2015unet}.
A stratified 15\% validation split (413 images) was held out from training for
model selection. Models were selected based on best validation Dice score,
and a maximum of 50 epochs was used per experiment.

\subsection{Architecture Choices}

Five experiments were designed to systematically ablate the effect of
architecture, initialisation strategy, loss function, and data augmentation.
Table~\ref{tab:experiments} summarises the configurations.

\begin{table}[H]
\centering
\caption{Summary of the five experimental configurations.}
\label{tab:experiments}
\begin{tabular}{llllll}
\toprule
\textbf{Exp} & \textbf{Architecture} & \textbf{Encoder} & \textbf{Init} & \textbf{Loss} & \textbf{Augmentation} \\
\midrule
1 & U-Net        & ResNet-34       & Random   & BCE                 & Light   \\
2 & U-Net        & ResNet-34       & ImageNet & BCE + Dice          & Moderate \\
3 & UNet++       & EfficientNet-B0 & ImageNet & Focal + Dice        & Heavy   \\
4 & FPN          & ResNet-50       & ImageNet & Dice                & Heavy   \\
5 & MAnet        & ResNet-34       & ImageNet & BCE + Dice + Focal  & Very Heavy \\
\bottomrule
\end{tabular}
\end{table}

\textbf{Experiment 1 -- Baseline (U-Net, Random Initialisation, BCE):}
This experiment serves as the controlled baseline. U-Net \cite{ronneberger2015unet} is the
seminal architecture for biomedical image segmentation, employing a symmetric
encoder-decoder structure with skip connections that preserve spatial detail.
Random initialisation was used to establish a lower-bound reference.
Binary cross-entropy (BCE) is the simplest pixel-level loss, treating segmentation
as independent binary classification for each pixel.
Minimal augmentation (horizontal and vertical flips) was applied to avoid confounding
the initialisation comparison.
The learning rate was set to $1 \times 10^{-3}$, which is standard for randomly
initialised networks.

\textbf{Experiment 2 -- U-Net with ImageNet Pretraining and Combined BCE+Dice Loss:}
This experiment isolates the benefit of transfer learning from ImageNet.
Pretraining on ImageNet provides rich feature representations -- edges, textures,
and hierarchical structures -- that transfer effectively even to grayscale medical images
\cite{tajbakhsh2016cnn}.
The BCE+Dice compound loss was used as it addresses the class imbalance inherent in
nerve segmentation (nerve pixels are a small fraction of the image).
BCE provides stable pixel-wise gradients while Dice directly maximises overlap,
making their combination a popular and well-supported choice in the literature
\cite{sudre2017generalised}.
Moderate augmentations (rotation $\pm15^\circ$, brightness/contrast jitter)
were added to reduce overfitting without distorting anatomical plausibility.
The learning rate was reduced to $1 \times 10^{-4}$ to protect pretrained weights.

\textbf{Experiment 3 -- UNet++ with EfficientNet-B0 and Focal+Dice Loss:}
UNet++ \cite{zhou2018unetplusplus} introduces dense skip connections with nested
encoder-decoder sub-networks, allowing the model to aggregate features at multiple
semantic scales. This design has shown improvements over standard U-Net on challenging
segmentation tasks by enabling flexible feature aggregation.
EfficientNet-B0 \cite{tan2019efficientnet} was chosen as the encoder due to its
favourable accuracy-to-parameter trade-off, achieved through compound scaling.
Focal loss \cite{lin2017focal} was combined with Dice to additionally penalise
misclassified hard examples -- particularly useful when nerve boundaries are
ambiguous or partially occluded.
Heavy augmentation (elastic deformation, grid distortion, CLAHE) was applied
to improve invariance to ultrasound-specific artefacts such as shadowing and
speckle noise.

\textbf{Experiment 4 -- FPN with ResNet-50 and Pure Dice Loss:}
Feature Pyramid Networks (FPN) \cite{lin2017fpn} build a top-down feature hierarchy
with lateral connections, enabling multi-scale detection. In segmentation, FPN decoders
have proven competitive with U-Net-style approaches by explicitly reasoning at multiple
resolutions, which is beneficial for nerves that vary considerably in appearance and size
across subjects.
ResNet-50 was selected as a deeper encoder variant to assess the effect of encoder
capacity.
Pure Dice loss was used to directly optimise the metric of interest, without the
regularisation effect of BCE.

\textbf{Experiment 5 -- MAnet with Combined Loss and Very Heavy Augmentation:}
Multi-scale Attention Net (MAnet) \cite{fan2020manet} employs channel and position
attention mechanisms at multiple scales to capture long-range dependencies, which may
be advantageous for localising the brachial plexus in the presence of surrounding
anatomical structures.
The combined loss (BCE + Dice + Focal, weighted 0.4:0.4:0.2) was intended to capture
the complementary supervision signals of all three objectives simultaneously.
Very heavy augmentation (GaussNoise, CoarseDropout, ShiftScaleRotate on top of all
previous augmentations) was employed as a form of regularisation given the model's
larger representational capacity.

\subsection{Common Hyperparameters}

All pretrained experiments used Adam \cite{kingma2014adam} with
$lr = 1 \times 10^{-4}$, $\beta_1 = 0.9$, $\beta_2 = 0.999$, weight decay $10^{-5}$.
Experiment 1 used $lr = 1 \times 10^{-3}$ due to random initialisation.
A ReduceLROnPlateau scheduler (factor=0.5, patience=5 epochs) monitored validation Dice
to adaptively reduce the learning rate upon plateaus.
Early stopping with patience of 12 epochs was used to prevent overfitting.
Batch size was 8 across all experiments.

%% ─────────────────────────────────────────────────────────────────────────────
\section{Task 1.2 -- Test Set Evaluation}
%% ─────────────────────────────────────────────────────────────────────────────

All five trained models were evaluated on the 2,878 test images.
Table~\ref{tab:results} reports Dice and Jaccard (primary metrics) alongside four
additional metrics: Precision, Recall, Accuracy, and 95th-percentile Hausdorff
Distance (HD95).

\begin{table}[H]
\centering
\caption{Test set performance across all five experiments (mean $\pm$ std).
HD95 is reported in pixels. Best result per metric shown in \textbf{bold}.}
\label{tab:results}
\footnotesize
\begin{tabular}{@{}l
  S[table-format=1.3]@{\,$\pm$\,}S[table-format=1.3]
  S[table-format=1.3]@{\,$\pm$\,}S[table-format=1.3]
  S[table-format=1.3]@{\,$\pm$\,}S[table-format=1.3]@{}}
\toprule
\textbf{Exp} & \multicolumn{2}{c}{\textbf{Dice}} & \multicolumn{2}{c}{\textbf{Jaccard}} & \multicolumn{2}{c}{\textbf{Precision}} \\
\midrule
1        & 0.706 & 0.410 & 0.676 & 0.416 & 0.882 & 0.292 \\
\textbf{2} & \textbf{0.744} & \textbf{0.396} & \textbf{0.717} & \textbf{0.399} & 0.902 & 0.254 \\
3        & 0.702 & 0.411 & 0.670 & 0.415 & 0.851 & 0.318 \\
4        & 0.732 & 0.402 & 0.705 & 0.406 & \textbf{0.906} & \textbf{0.249} \\
5        & 0.723 & 0.407 & 0.696 & 0.412 & \textbf{0.908} & \textbf{0.260} \\
\bottomrule
\end{tabular}

\vspace{6pt}

\begin{tabular}{@{}l
  S[table-format=1.3]@{\,$\pm$\,}S[table-format=1.3]
  S[table-format=1.3]@{\,$\pm$\,}S[table-format=1.3]
  S[table-format=1.2]@{\,$\pm$\,}S[table-format=1.3]@{}}
\toprule
\textbf{Exp} & \multicolumn{2}{c}{\textbf{Recall}} & \multicolumn{2}{c}{\textbf{Accuracy}} & \multicolumn{2}{c}{\textbf{HD95 $\downarrow$}} \\
\midrule
1        & 0.781 & 0.364 & 0.992 & 0.011 & 3.156 & 8.379 \\
\textbf{2} & \textbf{0.804} & \textbf{0.354} & \textbf{0.992} & \textbf{0.011} & \textbf{2.14} & \textbf{6.572} \\
3        & 0.802 & 0.346 & 0.992 & 0.011 & 3.303 & 9.124 \\
4        & 0.788 & 0.367 & 0.992 & 0.011 & 2.257 & 6.726 \\
5        & 0.777 & 0.370 & 0.992 & 0.011 & 2.848 & 8.509 \\
\bottomrule
\end{tabular}
\end{table}

\textbf{Experiment 2} achieves the highest Dice (0.744), Jaccard (0.717), Recall (0.804),
and the lowest HD95 (2.138), making it the top-performing model overall.
Experiments 4 and 5 achieve marginally higher Precision (0.906, 0.908) but at the
cost of lower Recall and overall Dice.
The high standard deviations across all experiments reflect the bimodal nature of
the dataset -- many images contain no nerve, leading to a mix of near-perfect
scores (true negatives) and failed detections.

Training curves for all experiments are shown in Figure~\ref{fig:curves}.

\begin{figure}[H]
    \centering
    \includegraphics[width=\textwidth]{figures3/training_curves.pdf}
    \caption{Validation loss and validation Dice score across training epochs for
    Experiments 1--4. The training history for Experiment 5 was not available
    as the job completed without saving epoch logs; its final checkpoint was
    preserved and used for test evaluation.}
    \label{fig:curves}
\end{figure}

%% ─────────────────────────────────────────────────────────────────────────────
\section{Task 1.3 -- Why Experiment 2 Outperforms the Others}
%% ─────────────────────────────────────────────────────────────────────────────

Experiment 2 (U-Net + ResNet-34 ImageNet + BCE+Dice, moderate augmentation)
achieved the best overall test performance. The following provides a scientific
analysis of the factors that drove this advantage.

\textbf{1. Transfer learning from ImageNet.}
The most significant factor distinguishing Exp2 from Exp1 is the use of
ImageNet-pretrained weights. The improvement in test Dice from 0.706 (random init)
to 0.744 (ImageNet pretrained) -- an absolute gain of 3.8\% -- is consistent with
the broad literature on transfer learning for medical imaging
\cite{tajbakhsh2016cnn, raghu2019transfusion}.
Pretrained features encode low-level (edges, gradients) to mid-level (textures,
shapes) representations that are directly applicable to detecting anatomical
structures in ultrasound, even though the source domain (natural images) differs
significantly. Random initialisation (Exp1) requires the model to learn all these
representations from scratch using only 2,344 training images, which is insufficient
for a deep encoder to converge to globally optimal features in 50 epochs.

\textbf{2. BCE + Dice loss versus alternatives.}
The BCE+Dice combination in Exp2 outperforms pure Dice (Exp4: 0.732) and Focal+Dice
(Exp3: 0.702).
Pure Dice loss (Exp4) lacks pixel-level regularisation and can be unstable on small,
spatially sparse targets -- it optimises a global ratio and provides weaker gradient
signal on individual hard pixels \cite{sudre2017generalised}.
Focal loss (Exp3) emphasises hard examples, but in this dataset the hardest cases are
images where no nerve is present at all -- focusing on these adds noise rather than
useful signal.
BCE provides stable per-pixel gradients throughout training, acting as a regulariser
that complements the region-level Dice term.
The combined BCE+Dice loss is therefore the best-balanced objective for this task.

\textbf{3. Appropriate architecture-augmentation balance.}
U-Net with a ResNet-34 encoder offers a strong capacity-to-data ratio for this
dataset size.
UNet++ (Exp3) and MAnet (Exp5) are architecturally more complex, requiring more data
to realise their theoretical advantages; in practice, they overfit or benefit less
from the available training set of approximately 2,300 images.
The evidence is in the early stopping behaviour: Exp3 stopped at epoch 32 and Exp4
at epoch 36, while Exp2 trained the full 50 epochs -- indicating that Exp2's design
made better use of all available training signal.
Very heavy augmentation (Exp5) over-regularises the learning signal for a ResNet-34
encoder, hurting the precision of gradient estimates and ultimately the model's
ability to detect the nerve reliably.
Moderate augmentation in Exp2 strikes the right balance -- sufficient to generalise
across patient variability without degrading the informativeness of the training
signal.

\textbf{4. Recall advantage.}
Exp2 achieves the highest Recall (0.804), which is clinically more important than
Precision in nerve detection: a false negative (missed nerve) is more dangerous than
a false positive (extra predicted region) in a surgical guidance context.
The BCE component of the loss ensures that nerve pixels are not ignored even when the
nerve occupies a small fraction of the image, thereby maintaining sensitivity.

%% ─────────────────────────────────────────────────────────────────────────────
\section{Task 1.4 -- Error Analysis on Worst 10\% of Cases}
%% ─────────────────────────────────────────────────────────────────────────────

\subsection{Identification of Failure Cases}

Using Experiment 2 (best model), the test samples were ranked by Dice score in
ascending order. The worst 10\% corresponds to 287 samples, all with Dice $\approx$ 0.
Figure~\ref{fig:errors} shows representative failure cases.

\begin{figure}[H]
    \centering
    \includegraphics[width=\textwidth]{figures3/task1_4_error_analysis.pdf}
    \caption{Worst 10\% failure cases from Experiment 2. Each row shows the original
    ultrasound image, the ground-truth mask, the predicted mask, and an overlay
    (ground truth in green, prediction in red).}
    \label{fig:errors}
\end{figure}

\subsection{Common Characteristics of Failure Cases}

Inspection of the failure cases reveals three dominant patterns:

\textbf{1. Absent or ambiguous nerve (empty ground truth).}
Approximately 394 test images ($\sim$13.7\%) have empty ground-truth masks,
meaning the nerve is either absent from the field of view or too faint to annotate.
The model tends to predict a small positive region in these cases, resulting in
Dice = 0 (false positive with no true positive).
This is the primary source of failure and is a fundamental challenge inherent to
the dataset.

\textbf{2. Highly variable nerve appearance.}
Ultrasound imaging is highly operator- and patient-dependent. In failure cases,
the nerve either (a) appears with very low contrast relative to surrounding tissue,
making it visually indistinguishable, or (b) is present but in an atypical position
or orientation due to patient posture variation.
The model has not encountered sufficient diversity in these presentations during training.

\textbf{3. Speckle noise and shadowing artefacts.}
Some failure cases show strong acoustic shadowing or high speckle noise that
obscures the nerve entirely. The model incorrectly segments the boundary of the
shadow or a nearby hypoechoic structure as the nerve.

\subsection{Suggested Improvements}

\textbf{1. Handle the empty-mask cases explicitly.}
A two-stage approach \cite{oktay2018attentionunet} -- a classifier to first predict
whether a nerve is present, followed by the segmentation model only when positive --
would eliminate false positives in empty-mask images and likely improve Dice
significantly.

\textbf{2. Contrastive or self-supervised pre-training on ultrasound data.}
Fine-tuning from ImageNet features is effective but limited by domain gap.
Pre-training the encoder on a large corpus of unlabelled ultrasound images using
contrastive learning (e.g., SimCLR \cite{chen2020simclr}) would yield
ultrasound-specific features more robust to speckle and artefacts.

\textbf{3. Test-time augmentation (TTA).}
Applying multiple augmentations (flips, rotations) at inference and averaging the
predictions is a simple but effective way to improve robustness on ambiguous cases
\cite{wang2019aleatoric}.
This adds no training cost and can improve Dice on low-contrast or boundary-ambiguous
cases where a single forward pass produces inconsistent results.

%% ─────────────────────────────────────────────────────────────────────────────
\section{Task 1.5 -- Qualitative Visualisation}
%% ─────────────────────────────────────────────────────────────────────────────

Figure~\ref{fig:qualitative} presents segmentation predictions from all five
experiments on five randomly selected test samples.
Each row corresponds to one sample; columns show the original image, the
ground-truth mask, and the predicted mask from each experiment.

\begin{figure}[H]
    \centering
    \includegraphics[width=\textwidth]{figures3/task1_5_qualitative.pdf}
    \caption{Qualitative segmentation results for five randomly selected test samples
    across all five experiments. Columns: original image, ground-truth mask,
    Exp1--Exp5 predictions (left to right).}
    \label{fig:qualitative}
\end{figure}

The qualitative results are consistent with the quantitative findings.
Experiment 2 produces the most precise delineation of the nerve boundary, with fewer
spurious activations in background regions.
Experiment 1 (random init) tends to produce noisier, less spatially precise masks,
reflecting the sub-optimal feature learning.
Experiments 3 and 5 occasionally over-segment, particularly in samples with
pronounced speckle noise, attributable to the focal loss component amplifying
ambiguous high-loss pixels.

%% ─────────────────────────────────────────────────────────────────────────────
\bibliographystyle{plain}
\begin{thebibliography}{99}

\bibitem{ronneberger2015unet}
O. Ronneberger, P. Fischer, and T. Brox,
``U-Net: Convolutional networks for biomedical image segmentation,''
in \textit{MICCAI}, 2015.

\bibitem{tajbakhsh2016cnn}
N. Tajbakhsh et al.,
``Convolutional neural networks for medical image analysis: Full training or fine tuning?,''
\textit{IEEE Trans. Med. Imaging}, vol. 35, no. 5, pp. 1299--1312, 2016.

\bibitem{raghu2019transfusion}
M. Raghu, C. Zhang, J. Kleinberg, and S. Bengio,
``Transfusion: Understanding transfer learning for medical imaging,''
in \textit{NeurIPS}, 2019.

\bibitem{zhou2018unetplusplus}
Z. Zhou, M. M. R. Siddiquee, N. Tajbakhsh, and J. Liang,
``UNet++: A nested U-Net architecture for medical image segmentation,''
in \textit{MICCAI Deep Learning Workshop}, 2018.

\bibitem{tan2019efficientnet}
M. Tan and Q. V. Le,
``EfficientNet: Rethinking model scaling for convolutional neural networks,''
in \textit{ICML}, 2019.

\bibitem{lin2017focal}
T.-Y. Lin, P. Goyal, R. Girshick, K. He, and P. Dollár,
``Focal loss for dense object detection,''
in \textit{ICCV}, 2017.

\bibitem{lin2017fpn}
T.-Y. Lin, P. Dollár, R. Girshick, K. He, B. Hariharan, and S. Belongie,
``Feature pyramid networks for object detection,''
in \textit{CVPR}, 2017.

\bibitem{fan2020manet}
T. Fan, G. Wang, Y. Li, and H. Wang,
``MA-Net: A multi-scale attention network for liver and tumor segmentation,''
\textit{IEEE Access}, vol. 8, pp. 179656--179665, 2020.

\bibitem{sudre2017generalised}
C. H. Sudre, W. Li, T. Vercauteren, S. Ourselin, and M. J. Cardoso,
``Generalised Dice overlap as a deep learning loss function for highly unbalanced segmentations,''
in \textit{MICCAI DLMIA Workshop}, 2017.

\bibitem{kingma2014adam}
D. P. Kingma and J. Ba,
``Adam: A method for stochastic optimization,''
in \textit{ICLR}, 2015.

\bibitem{oktay2018attentionunet}
O. Oktay et al.,
``Attention U-Net: Learning where to look for the pancreas,''
in \textit{MIDL}, 2018.

\bibitem{chen2020simclr}
T. Chen, S. Kornblith, M. Norouzi, and G. Hinton,
``A simple framework for contrastive learning of visual representations,''
in \textit{ICML}, 2020.

\bibitem{wang2019aleatoric}
G. Wang, W. Li, M. Aertsen, J. Deprest, S. Ourselin, and T. Vercauteren,
``Aleatoric uncertainty estimation with test-time augmentation for medical image segmentation,''
\textit{Neurocomputing}, vol. 338, pp. 34--45, 2019.

\end{thebibliography}

\end{document}
